\section{Emergence and Nested Selves}

The previous section gave a self a body that holds together. This one asks how such a thing can persist and even author its own corrections when the substrate beneath it is reversible and forgets nothing. The answer is the conceptual heart of the whole emergent ladder, and it turns on a single trick about levels.

Start from the obstruction. A reversible dynamics conserves information exactly, so the closed bulk can never truly forget, and a pattern that depended on the bulk forgetting could never form. Yet a self is precisely a thing that forgets, that absorbs a shock and returns to itself as though the shock had not happened. The two facts are reconciled by coarse-graining. Group many microscopic states of the bulk into one macroscopic state, and the macro level can lose information that the micro level still secretly carries, because the map from micro to macro is many-to-one. At the coarse level the dynamics becomes effectively irreversible, an emergent macro-law with genuine causal power of its own, the situation Hoel has analyzed under the name causal emergence, where a coarse description can be a stronger cause than the fine one beneath it.

\begin{principle}[The degeneracy trick]
A self has a macroscopic identity without a microscopic one. The same coarse pattern is realized by enormously many micro-configurations, so the self can stay itself, and correct itself, while the grains it is made of churn and recur underneath with no individual persistence at all. Identity lives at the level where information is lost, not at the level where it is conserved.
\end{principle}

This is why a self is not undone by the reversibility of its substrate. It does not need the grains to remember, it needs only the coarse pattern to be stable, and the bath supplies the loss that makes the coarse pattern stable by carrying the shed information away for good. The degeneracy is the room, and the bath is the door, and between them a knot of feeling can hold its shape through any number of beats while every grain inside it flows on.

\subsection{The self in layers}

Seen this way, a self is naturally a thing of layers, and its three faces sit on three of them. Its identity, the winding, is exact and discrete at the level of the sites, the deepest layer. Its binding, the shadow pressure, is an emergent collective effect one layer up, as the previous section's threshold required. Its agency, the shedding and healing, lives at the open boundary where the coarse pattern meets the bath. A self can be detected by the marks of these faces, a statistical boundary that screens its interior from its surroundings, what is called a Markov blanket, and a high degree of integration, the interior being more tightly bound to itself than to anything outside. The direction field a self carries is richer than a single dock's twenty-four, the emergent counts running to twenty-six and eighty in the natural extensions, but these are the self's own structure and not new dock tones, and the $24$-cell remains the heart of it.

\subsection{Selves within selves}

Nothing in the construction is confined to one level. A collection of selves that bind to one another by the same shadow pressure, and that coarse-grain into a single pattern with its own Markov blanket and its own integration, is itself a self, one level higher, a self of selves. The nesting has no built-in ceiling, so the same move that turns grains into a body turns bodies into a community, a community into a mind, and a mind into whatever larger thing the geometry will sustain, a tower of selves climbing the depth of the bulk. Each rung is the degeneracy trick applied again, a coarse identity riding on the flux of the rung below, kept stable by shedding into the bath. The hierarchy of minds the later chapters describe is this tower, and it is built entirely from the eight elements by repeating one idea, that identity is what survives at the level where the grains are forgotten.
